\documentclass[12pt]{article}
\usepackage{amsmath, amssymb, graphicx, float}
\usepackage[shortlabels]{enumitem}
\usepackage{geometry}
\usepackage{pgfplots}
\pgfplotsset{compat=1.18}

\geometry{margin=1in}

\begin{document}

\begin{center}
\textbf{ELEC 327 Lab \#6} \\[6pt]
Liam Brady \\
Due: Saturday 3/7/26
\end{center}

The music is all logic that came from Lab 5 (expect for the fact that the buzzer is now driven by the DAC), so nothing changed in the state logic. Therefore it was very easy to hook up both song and button led reactions to this already defined logic.
\indent For each different note in the song, a color is defined. Then, when that note is being played, all leds turn on, all of them being inverted except for one, which moves around in a circle. The speed it moves is equal to the length of an eighth note, so in a quarter note the led will move from the first led to the second, and for a whole note the led will move around the 4 leds twice. In the pause between notes, all leds will be turned off.
\indent For the buttons, each button is also related to a color, set equivalently to the colors of Simon Says. Then, when that button is pressed, that specific color is sent.
\end{document}
